\documentclass[11pt]{article}
\usepackage[a4paper,margin=1in]{geometry}
\usepackage[T1]{fontenc}
\usepackage{amsmath,booktabs,array,tabularx}
\usepackage[authoryear,round]{natbib}
\usepackage[hidelinks]{hyperref}
\newcolumntype{Y}{>{\raggedright\arraybackslash}X}
\newcommand{\Neff}{N_{\mathrm{eff}}}
\title{Patient Shortage in Histopathology:\\Does Cohort Composition Matter\\at a Fixed Patient Count?}
\author{Yannik Queisler}
\date{September 13, 2026}

\begin{document}
\raggedbottom
\widowpenalty=10000
\clubpenalty=10000
\setlength{\emergencystretch}{2em}
\maketitle
\begin{abstract}
\noindent TCGA-UT patch classifiers trained on more patients are more accurate than redundancy-adjusted effective support predicts, and tissue source sites explain little of this residual. Two tests of patient coverage stopped at their manipulation checks, because a twenty-patient allocation nested around five seed patients covers held-out patients almost as well as a random one, whichever patients are added. This experiment therefore holds the patient count fixed and varies only which patients form the cohort. Three allocations of twenty patients with eight patches each share a randomly drawn focal patient: its nineteen nearest neighbours in the patient-mean Virchow2 embedding (\emph{clustered}), nineteen randomly drawn patients (\emph{random}), and nineteen patients chosen by facility location on the training pool (\emph{dispersed}). Patients, patches, and effective support are equal across the three allocations. The primary contrast, \emph{concentration damage}, is the accuracy of the random minus the clustered allocation; the secondary contrast, \emph{selection gain}, is the accuracy of the dispersed minus the random allocation. A manipulation check, run before any classifier is trained, requires the clustered-to-random coverage change to reach at least half of the coverage change from five to twenty random patients in every split. Embeddings from a preceding diagnostic give 0.58 to 0.68 for this ratio but only 0.26 to 0.27 for the random-to-dispersed change, so the selection gain is interpreted only as a gain, never as evidence of equivalence. Concentration damage above one percentage point would show that cohort composition matters at a fixed patient count; an additional selection gain above one point would support coverage-guided patient selection; practical equivalence for concentration damage would weaken patient-mean coverage as the missing source of information.
\end{abstract}
\setcounter{tocdepth}{1}
\tableofcontents
\clearpage

\section{Motivation}
\label{sec:motivation}
Spreading a labelling budget over more patients improves TCGA-UT patch classification beyond what redundancy explains. At 160 patches per class, twenty patients contributing eight patches each outperform five patients contributing 32 patches each by 9.52 percentage points of patient-macro balanced accuracy \citep{queisler2026effectiveSupport}. At equal effective support, each doubling of patients still adds 2.62 points \citep{queisler2026multidirectionalRedundancy}, about 5.2 points over the two doublings of this comparison. Doubling tissue source sites at a fixed patient count adds only 0.84 points \citep{queisler2026siteCoverage}.

\emph{Patient coverage}, the share of a cancer type's between-patient variation that the training patients represent, remains the leading label-free candidate. Its two tests compared a low-coverage and a random twenty-patient allocation that both contained the same five seed patients \citep{queisler2026patientCoverage,queisler2026concentratedCoverage}. Both stopped at the manipulation check: the low-coverage allocation attained 84 to 88 percent, and after concentration around a single seed still 67 to 73 percent, of the random allocation's coverage improvement. A diagnostic of seven nested rules reached no leakage below 0.53 \citep{queisler2026concentratedCoverage}. Quadrupling the patient count shortens the distance from most held-out patients to their nearest training patient, whichever patients are added, so a step that raises patient count at nearly constant coverage is out of reach for this coverage measure.

The same diagnostic shows that coverage varies strongly between cohorts of equal size. Twenty patients forming one neighbourhood cover validation patients at a mean distance of 0.346 to 0.356, twenty random patients at 0.292 to 0.296, and twenty patients selected by facility location at 0.268 to 0.271; five random patients cover them at 0.385 to 0.387. This report therefore replaces the question whether similar patients help beyond effective support with a question that the measure can address: \emph{at a fixed patient count, does a cohort that covers the class's patients poorly train a worse classifier, and does one that covers them well train a better classifier than random sampling?}

\section{Study design}
\label{sec:design}
Only the composition of a twenty-patient cohort changes. The design inherits from \citet{queisler2026concentratedCoverage} the 30 eligible TCGA-UT cancer types, the frozen 2,560-dimensional Virchow2 features \citep{zimmermann2024virchow2}, the three fixed patient-disjoint splits with their natural validation and test partitions, the nested round-robin patch order, and the readout: multinomial logistic regression with its regularization strength selected on the validation partition for every fit. It also inherits the patient distance. For split $r$ and class $c$, $\bar{\mathbf x}_{irc}$ is the mean feature vector of patient $i$'s patches of the class, $\boldsymbol\mu_r$ the average of these means over all classes, and
\begin{equation}
 \mathbf e_{irc}=\frac{\bar{\mathbf x}_{irc}-\boldsymbol\mu_r}{\lVert\bar{\mathbf x}_{irc}-\boldsymbol\mu_r\rVert},
 \qquad
 d(i,j)=1-\mathbf e_{irc}^{\top}\mathbf e_{jrc}.
 \label{eq:distance}
\end{equation}

\begin{table}[htbp]
\centering
\renewcommand{\arraystretch}{1.15}
\begin{tabular}{@{}lrrl@{}}
\toprule
Allocation & Patients $G$ & Patches per patient $m$ & Patients \\
\midrule
Clustered & 20 & 8 & focal patient, its 19 nearest neighbours \\
Random & 20 & 8 & focal patient, 19 random patients \\
Dispersed & 20 & 8 & focal patient, 19 facility-location patients \\
\bottomrule
\end{tabular}
\caption{The three allocations share the focal patient, the patient count, the budget of 160 patches per class, and therefore the effective support. They differ only in how the remaining nineteen patients are chosen, which moves coverage from low to high.}
\label{tab:allocations}
\end{table}

\noindent Table~\ref{tab:allocations} lists the allocations. Because all three share $G=20$ and $m=8$, effective support is identical by construction, and no effective-support surface enters the analysis.

\subsection{Allocation rules}
\label{sec:allocation}
Three rules define the training sets and check the manipulation before any classifier is trained. First, every class is admitted. The pool $\mathcal E_{rc}$ contains the training patients with at least 32 patches of class $c$; pools hold between 30 and 514 patients, so every class can supply twenty.

Second, the allocations share a draw structure. For each split, class, and draw $d\in\{1,\ldots,5\}$, the five \emph{seed} patients are those of the stored five-patient draw of \citet{queisler2026effectiveSupport}, and the \emph{focal patient} $f$ is the seed drawn uniformly at random in \citet{queisler2026concentratedCoverage}. The random allocation is the random allocation of \citet{queisler2026concentratedCoverage}: the five seeds and fifteen randomly drawn pool patients. It contains $f$ and is distributed like a twenty-patient random draw. The clustered allocation contains $f$ and the nineteen patients of $\mathcal E_{rc}\setminus\{f\}$ closest to it under Equation~\eqref{eq:distance}. The dispersed allocation starts from $S_0=\{f\}$ and adds nineteen patients greedily by \emph{facility location},
\begin{equation}
 S_{k+1}=S_k\cup\Bigl\{\operatorname*{arg\,min}_{j\in\mathcal E_{rc}\setminus S_k}\ \frac{1}{|\mathcal E_{rc}|}\sum_{i\in\mathcal E_{rc}}\ \min_{l\in S_k\cup\{j\}} d(i,l)\Bigr\},
 \label{eq:facility}
\end{equation}
\noindent so each step adds the pool patient that most shortens the mean distance from pool patients to their nearest selected patient. Ties in both rules are broken by patient identifier, and neither rule uses validation or test patients. Every patient contributes the first eight patches of its round-robin order. The design requires 45 fits, one per allocation, split, and draw.

Third, a manipulation check compares how well the allocations cover held-out patients. The \emph{coverage distance} of a patient set is the mean distance from the class's validation patients $\mathcal V_{rc}$ to their nearest training patient,
\begin{equation}
 r_{rc}(\mathcal P)=\frac{1}{|\mathcal V_{rc}|}\sum_{i\in\mathcal V_{rc}}\min_{j\in\mathcal P}d(i,j).
 \label{eq:coverage}
\end{equation}
\noindent With $r_{\mathrm{clu}}$, $r_{\mathrm{ran}}$, and $r_{\mathrm{dis}}$ the coverage distances of the three allocations and $r_{5}$ that of the five seeds, each averaged over classes and draws within a split, the coverage change from five to twenty random patients serves as the reference scale. The two \emph{coverage ratios}
\begin{equation}
 \rho_{\mathrm{con}}=\frac{r_{\mathrm{clu}}-r_{\mathrm{ran}}}{r_{5}-r_{\mathrm{ran}}},
 \qquad
 \rho_{\mathrm{sel}}=\frac{r_{\mathrm{ran}}-r_{\mathrm{dis}}}{r_{5}-r_{\mathrm{ran}}}
 \label{eq:ratios}
\end{equation}
\noindent express the coverage change behind each contrast in units of this scale. Fitting proceeds only if $\rho_{\mathrm{con}}\ge 0.5$ in all three splits. The ratio $\rho_{\mathrm{sel}}$, the class-level ratios, the mean pairwise distance within each cohort, and the share of patients from the focal patient's tissue source site are recorded without a threshold.

Table~\ref{tab:preview} reports these quantities from the diagnostic of \citet{queisler2026concentratedCoverage}, which used the same embeddings, draws, and selection rules. The check is expected to pass: $\rho_{\mathrm{con}}$ lies between 0.58 and 0.68, and the class-level ratio reaches 0.5 in 52 of 90 class--split combinations. The secondary contrast rests on a much smaller manipulation, a quarter of the reference scale, which constrains its interpretation in Section~\ref{sec:interpretation}.

\begin{table}[htbp]
\centering
\renewcommand{\arraystretch}{1.15}
\begin{tabular}{@{}lrrrrrr@{}}
\toprule
 & \multicolumn{4}{c}{Coverage distance} & & \\
\cmidrule(lr){2-5}
Split & $r_{5}$ & $r_{\mathrm{clu}}$ & $r_{\mathrm{ran}}$ & $r_{\mathrm{dis}}$ & $\rho_{\mathrm{con}}$ & $\rho_{\mathrm{sel}}$ \\
\midrule
1 & 0.387 & 0.353 & 0.296 & 0.271 & 0.63 & 0.27 \\
2 & 0.385 & 0.356 & 0.293 & 0.268 & 0.68 & 0.27 \\
3 & 0.385 & 0.346 & 0.292 & 0.268 & 0.58 & 0.26 \\
\bottomrule
\end{tabular}
\caption{Clustering moves coverage far more than facility location does. Coverage distances are mean cosine distances from validation patients to their nearest training patient, averaged over the 30 classes and five draws, from the diagnostic of \citet{queisler2026concentratedCoverage}. Mean pairwise distances within a cohort are 0.35 to 0.37 for clustered, 0.59 to 0.60 for random, and 0.67 for dispersed allocations.}
\label{tab:preview}
\end{table}

\section{Evaluation}
\label{sec:evaluation}
Both contrasts are read with the same accuracy measure and the same uncertainty procedure.

\subsection{Concentration damage and selection gain}
\label{sec:contrasts}
Accuracy is patient-macro balanced accuracy over all 30 classes: patch recall is averaged within each test patient and true class, then across patients within each class, and finally across classes. Averaging three splits and five draws gives $A_{\mathrm{clu}}$, $A_{\mathrm{ran}}$, and $A_{\mathrm{dis}}$ in percent. The primary endpoint is the \emph{concentration damage}
\begin{equation}
 \Delta_{\mathrm{con}}=A_{\mathrm{ran}}-A_{\mathrm{clu}},
 \label{eq:damage}
\end{equation}
\noindent the accuracy lost when nineteen patients resemble one patient instead of being drawn at random. The secondary endpoint is the \emph{selection gain}
\begin{equation}
 \Delta_{\mathrm{sel}}=A_{\mathrm{dis}}-A_{\mathrm{ran}},
 \label{eq:gain}
\end{equation}
\noindent the accuracy won when coverage-maximizing selection replaces random sampling. Both hold patients, patches per patient, budget, and effective support fixed.

A rough scale helps to read $\Delta_{\mathrm{con}}$. If coverage alone carried the residual of about 5.2 points between five and twenty random patients, and accuracy changed linearly with coverage distance, the concentration damage would be about $\rho_{\mathrm{con}}\times 5.2\approx 3$ points. This scale is descriptive, assumes a linearity that no preceding experiment tested, and carries no decision threshold.

Two secondary analyses describe where a concentration damage arises. First, each test patient receives per draw its \emph{coverage improvement}, the distance to its nearest clustered training patient minus the distance to its nearest random training patient, and $\Delta_{\mathrm{con}}$ is computed within tertiles of this improvement with the patient weights of the primary endpoint. Patient coverage predicts a damage that grows with the improvement. Second, all quantities, including the class-level coverage ratios, are reported separately for the thirteen site classes of \citet{queisler2026siteCoverage} and the seventeen remaining classes. These analyses carry no decision threshold.

\subsection{Uncertainty}
\label{sec:uncertainty}
Uncertainty follows \citet{queisler2026effectiveSupport}. Each of 2,000 paired Bayesian bootstrap replicates assigns $\mathrm{Dirichlet}(1,\ldots,1)$ weights to the distinct test patients \citep{rubin1981bayesian}, shared across allocations, splits, draws, and classes. The 2.5th and 97.5th percentiles give exploratory 95\% intervals. Embeddings and allocations are held fixed. The standard deviation of each allocation's accuracy and of both contrasts across the fifteen split--draw combinations is reported as a descriptive measure of sensitivity to the drawn patients.

\section{Interpretation criteria}
\label{sec:interpretation}
Conclusions rest on the test-patient intervals for $\Delta_{\mathrm{con}}$ and $\Delta_{\mathrm{sel}}$ and on the practical threshold of one percentage point used in the preceding experiments (Table~\ref{tab:readings}). The two contrasts are not treated symmetrically. Concentration damage rests on a coverage change of more than half the reference scale, so an interval within $[-1,1]$ counts as practical equivalence. Selection gain rests on a quarter of that scale, so only an interval above one point is informative; an interval within $[-1,1]$ does not show that coverage-guided selection is useless.

\begin{table}[htbp]
\centering
\renewcommand{\arraystretch}{1.15}
\begin{tabularx}{\textwidth}{@{}lYYY@{}}
\toprule
Reading & Claim & Concentration damage $\Delta_{\mathrm{con}}$ & Selection gain $\Delta_{\mathrm{sel}}$ \\
\midrule
Selection beyond sampling & Composition matters, and coverage-guided selection improves on random sampling & Interval above 1 & Interval above 1 \\
\addlinespace
Concentration damage & Composition matters, but only a poorly covering cohort is shown to hurt & Interval above 1 & Interval not above 1 \\
\addlinespace
No composition effect & Patient-mean coverage does not affect accuracy at twenty patients & Interval within $[-1,1]$ & Interval not above 1 \\
\bottomrule
\end{tabularx}
\caption{Each reading predicts a distinct combination of the two contrasts. Intervals are test-patient 95\% intervals in percentage points. All other combinations, including a concentration-damage interval that crosses a threshold, are inconclusive.}
\label{tab:readings}
\end{table}

\paragraph{Selection beyond sampling.}
If both contrasts exceed the threshold, patient-mean coverage separates useful from less useful cohorts in both directions from random sampling. Coverage distance would then be a candidate shortage signal alongside effective support, and the next step would compare coverage-selected and random cohorts of five patients with 32 patches each, the setting in which patient shortage was first observed.

\paragraph{Concentration damage.}
If only the concentration damage exceeds the threshold, a cohort that misses most of the class's between-patient variation trains a worse classifier, but random sampling may already capture most of the attainable coverage at twenty patients. Coverage would remain a plausible explanation of the breadth benefit; its value for patient selection would still be open and would again be tested first at five patients, where random cohorts cover the class less well.

\paragraph{No composition effect.}
If the concentration damage is practically zero, a coverage change of more than half the reference scale leaves accuracy unchanged, and patient-mean coverage is unlikely to carry the residual breadth benefit. The missing information would then lie in properties that a patient mean does not express, and the next candidate would be a distance between the patch distributions of patients.

\section{Limitations}
\label{sec:limitations}
The clustered allocation does more than withhold coverage. Its patients resemble each other, so they are more redundant with one another than effective support, which assumes independent patients, counts them; and they over-represent one appearance of the cancer type, towards which a classifier may shift its class boundary. A concentration damage therefore reflects between-patient redundancy, within-class imbalance of appearances, and missing coverage together, and the coverage-improvement tertiles separate these only partly.

Neighbours of one patient may share its tissue source site, because site signatures are visible in deep features of TCGA slides \citep{howard2021impact}. In the diagnostic, 24 to 27 percent of clustered and 16 to 18 percent of dispersed patients share the focal patient's site, and the site gain of 0.84 points \citep{queisler2026siteCoverage} gives a reference scale for how much of a concentration damage sites alone could explain.

The manipulation is uneven across cancer types. In small or homogeneous classes such as adrenocortical carcinoma, pancreatic adenocarcinoma, uterine carcinosarcoma, and uveal melanoma, the class-level ratio of the concentration contrast stays below 0.25 in most splits, so these classes contribute little contrast to the class average. Coverage is defined on patient means in one representation and a single cohort size, so a null result excludes coverage only as measured here. All results are conditional on frozen Virchow2 features, multinomial logistic regression, and the three overlapping splits, which are not independent replications.

\bibliographystyle{plainnat}
\bibliography{references}
\end{document}
